%% AIxB 2026 -- Camera-ready (Short Paper, 4 pages)
%% Submission 18: "A Graph Neural Network for Explainable Anti-Money
%% Laundering in Blockchain Transactions"
%%
%% REVISION NOTES (see RESPONSE_TO_REVIEWERS.md for the full point-by-point):
%%   R1/R4  code availability      -> Sec. I footnote + Sec. V-A
%%   R2/R4  multi-seed + stats     -> all tables now mean +/- std over 5 seeds,
%%                                    paired t-tests with 95% CIs
%%   R3     abstract / clarity     -> abstract rewritten, abbreviations expanded,
%%                                    new plain-language framing in Sec. I
%%   R4     temporal claim         -> scoped to timestep position encoding;
%%                                    TGN/TGAT explicitly listed as untested
%%   R4     component ablation     -> new Table IV (2x2 design)
%%   R4     reference [10]         -> resolved to Johannessen & Jullum (2025)

\documentclass[conference]{IEEEtran}
\IEEEoverridecommandlockouts

\usepackage{cite}
\usepackage{amsmath,amssymb}
\usepackage{graphicx}
\usepackage{booktabs}
\usepackage{url}
\usepackage[hidelinks]{hyperref}

\begin{document}

\title{A Graph Neural Network for Explainable\\
Anti-Money Laundering in Blockchain Transactions}

\author{\IEEEauthorblockN{Nandita Menon, Anushree Ashok, Jillian Priscilla, Debashis Guha}
\IEEEauthorblockA{Department of Artificial Intelligence\\
S P Jain School of Global Management, Dubai\\
\{nandita.as24dxb015, anushree.as24dxb017, jillian.as24dxb012\}@spjain.org, debashis.guha@spjain.org}}

\maketitle

\begin{abstract}
Criminals can launder Bitcoin by splitting illicit funds across multiple
transactions and wallet addresses, making individual payments difficult to
identify as suspicious. For compliance teams, the challenge is not only to
detect potentially illicit transactions but also to provide clear, defensible
reasons for investigating them. We formulate Bitcoin anti-money laundering
(AML) as a graph-based transaction classification problem and evaluate graph
neural networks (GNNs), which leverage both transaction features and
information from neighbouring transactions, against conventional classifiers
that consider each transaction independently. We introduce RiskMAGNN, a
heterogeneous GNN that explicitly models the wallet address connecting pairs
of transactions and weights each wallet-sharing relationship using priors
informed by financial-crime investigation. Across ten random seeds with
paired significance tests, graph structure provides the largest improvement,
increasing average precision by 37.6 percentage points over a feature-only
baseline ($p<10^{-4}$); RiskMAGNN further improves the area under the ROC
curve by 2.5 points over a standard heterogeneous GNN ($p=0.002$), although
its average-precision gain is not statistically significant across seeds.
Ensembling predictions across the ten seeds produces the strongest detector
(71.2\% average precision), and occasional training collapses in individual
runs highlight the limitations of single-seed results on this benchmark. We
further show that one of the three wallet-sharing relation types is entirely
absent from the training portion of the standard benchmark split; this
dataset issue, rather than model capacity, accounts for most of the observed
validation--test gap. An explanation module provides transaction-level
justifications to support compliance review, and all code and experimental
results are released for reproducibility.
\end{abstract}

\begin{IEEEkeywords}
anti-money laundering, graph neural networks, heterogeneous graphs, metapath
aggregation, blockchain, explainability
\end{IEEEkeywords}

\section{Introduction}

Money laundering hides the origin of criminal funds through chains of
intermediate transfers---\emph{layering}. On Bitcoin an operator splits a
payment across many transactions and wallet addresses, so each transfer looks
ordinary while the chain moves value away from its source. Two properties
shape any detection system. First, the evidence is relational: what makes a
transfer suspicious is the pattern of transfers around it, not its own amount
or fee. Second, regulation (the EU Anti-Money Laundering Directive, the US
Bank Secrecy Act) requires that a flagged transaction carry a
human-interpretable justification, so an accurate but opaque detector is not
deployable.

Graph neural networks (GNNs) fit the first property: a GNN classifies each
transaction by repeatedly mixing in information from its neighbours in the
fund-flow graph, capturing multi-hop structure a per-transaction classifier
cannot see~\cite{ferretti2025}. Most published benchmarks, however, model
transactions only. Real laundering schemes reuse wallet addresses, and a
transaction-only graph cannot express that two payments touched the same
wallet; \emph{heterogeneous} datasets that model addresses as first-class
nodes make this structure available~\cite{song2023}.

This paper reports a careful empirical study rather than a large architectural
gain. Our contributions: (i) a multi-seed comparison of feature-only,
homogeneous-graph, and heterogeneous-graph models on the Elliptic and HBTBD
datasets, with paired significance tests; (ii) RiskMAGNN, a domain-adapted
variant of the metapath-aggregated GNN (MAGNN)~\cite{fu2020} combining a
translation-based encoding of the shared wallet address~\cite{bordes2013}
with attention over relation types initialised from forensic priors, with a
component ablation; (iii) evidence that seed-to-seed spread on this benchmark
is comparable to the differences between architectures, so single-seed
comparisons should be treated with caution; and (iv) a diagnostic showing
that one of the three wallet-sharing relation types is entirely missing from
the official HBTBD training split---a previously unreported dataset issue
that accounts for most of the validation--test gap. An explanation module
built on GNNExplainer~\cite{ying2019} supplies per-transaction
justifications. Code and all result files are available at
\url{https://github.com/Iroh02/GNN_AML}.\footnote{See
\texttt{REPRODUCIBILITY.md} there for the exact command behind every number
in this paper, with expected runtimes.}

\section{Related Work}

Weber et al.~\cite{weber2019} established Elliptic as a public benchmark and
showed that graph convolutional networks beat feature-only classifiers;
Ferretti et al.~\cite{ferretti2025} confirmed consistent gains from graph
structure across GCN, GAT and GraphSAGE~\cite{hamilton2017}, while
identifying a ceiling for static homogeneous models on Elliptic's
chronological split. Lawal~\cite{lawal2024} integrated GNNExplainer and SHAP
into a homogeneous GCN, showing interpretability and performance are
compatible.

Deprez et al.~\cite{deprez2024} argue that temporal graph learning is needed
to model how layering unfolds, and architectures such as TGN and TGAT target
dynamic financial graphs. We do not test those architectures; we test only
whether supplying the timestep as an input feature helps, find the effect
inconsistent and not significant (Section~\ref{sec:temporal}), and therefore
neither support nor contradict that position.

Johannessen and Jullum~\cite{johannessen2025} show heterogeneous GNNs with
relation-specific message passing outperform homogeneous models on
proprietary bank data. Song and Gu~\cite{song2023} operationalise this for
Bitcoin with HBTBD, reaching 92.2\% precision and 73.9\% F1 by feeding
MAGNN-extracted features into a separately trained Random Forest. RiskMAGNN
instead learns end to end and injects domain knowledge directly into the
attention over relation types.

\section{Datasets}

\textbf{Elliptic} contains 203{,}769 transaction nodes and 234{,}355 directed
edges across 49 time steps, with 165 features per node (94 local properties;
71 neighbourhood aggregates) plus the time step; a subset is labelled licit
or illicit. We use the standard chronological split: train on steps 1--29
(10.88\% illicit), validate on 30--39, test on 40--49 (5.69\% illicit)---a
realistic distribution shift.

\textbf{HBTBD}~\cite{song2023} models wallet addresses as explicit nodes.
Three address roles---input (Add\_A), output/forwarding (Add\_B), change
(Add\_C)---define three \emph{metapaths} connecting transactions through a
shared address: \textbf{M1} (Tx\,$\rightarrow$\,Add\_A\,$\leftarrow$\,Tx),
both paying into the same address, the signature of a mixing service;
\textbf{M2} (Tx\,$\rightarrow$\,Add\_B\,$\rightarrow$\,Tx), one-way flow
through a forwarding address, characteristic of layering; \textbf{M3}
(Tx\,$\leftarrow$\,Add\_C\,$\rightarrow$\,Tx), both drawing from the same
source, consistent with structuring. The official split provides 29{,}699
training and 16{,}346 test transaction nodes with 165 features each. One
asymmetry is critical to every result below: \textbf{M2 edges are entirely
absent from the training set} (0 instances) but number 77{,}418 in the test
set, because one-way flows appear only in time steps 35--49. Any temporally
valid partition of HBTBD inherits this asymmetry---a severe distribution
shift for precisely the metapath that encodes layering.

\section{RiskMAGNN}

RiskMAGNN adapts MAGNN~\cite{fu2020} through two modifications.

\textbf{Translation-based metapath encoding.} Standard MAGNN aggregates a
metapath instance by concatenating its endpoint embeddings, discarding the
intermediate entity. For AML the wallet address in the middle carries
independent forensic signal: an address that recurs between illicit
transactions is itself evidence. We replace concatenation with a TransE-style
encoding~\cite{bordes2013}: node features are projected to a shared
$d$-dimensional space $h'$, and for instance
$P(\text{Tx}_v,\text{Add},\text{Tx}_u)$:
\begin{equation}
t = h'_{\text{Add}},\quad
o_1 = h'_{\text{Tx}_v} + t \odot r_1,\quad
o_2 = h'_{\text{Tx}_u} + t \odot r_2
\end{equation}
\begin{equation}
h_P = o_1 \oplus t \oplus o_2
\end{equation}
where $r_1, r_2 \in \mathbb{R}^d$ are learned relation vectors, $\odot$ is the
Hadamard product and $\oplus$ is concatenation. Because HBTBD ships
pre-aggregated transaction-to-transaction adjacencies rather than address
features, we approximate $t$ by mean-pooling the metapath neighbours of $v$.

\textbf{Risk-biased attention over relation types.} After mean-pooling
instances within each metapath type we fuse the three embeddings with
attention whose logits carry a learnable per-type bias $b_m$:
\begin{equation}
\alpha_m = \mathrm{softmax}_m\!\left(\frac{q^\top k_m}{\sqrt{d}} + b_m\right),
\qquad
h_v = \sum_{m} \alpha_m h_v^{(m)}
\end{equation}
with $q = W_q h_v$ and $k_m = W_k h_v^{(m)}$. The biases are
\emph{initialised} at $b_{\text{M2}}=0.6 > b_{\text{M1}}=0.3 >
b_{\text{M3}}=0.1$ (one-way flow above address reuse above shared sourcing)
and then trained. The model does not follow the prior blindly: the initial
biases alone would give M2 an attention share of $0.43$, but the learned mean
M2 attention is $0.29$ on the official split (M2 training edges absent) and
$0.17$ on the resplit of Section~\ref{sec:resplit} (M2 observed but
sparse)---the model down-weights the metapath the data cannot support.

Each layer applies a residual connection followed by LayerNorm; we stack
$L\in\{2,3\}$ layers with $d\in\{128,192\}$ and a two-layer classifier.
Training uses focal loss ($\alpha=0.75$, $\gamma=2.0$), AdamW
($\text{lr}=10^{-3}$, weight decay $5\times10^{-4}$) and a warmup-plus-cosine
schedule.

\textbf{Explanations.} We adapt GNNExplainer~\cite{ying2019} to the
heterogeneous setting by learning one soft edge mask $M_e^{(m)}$ per metapath
type plus a global feature mask $M_f$, optimised to preserve the original
prediction under an $\ell_1$ sparsity penalty, so explanations decompose
neighbour influence by relationship type.

\section{Experimental Results}

\subsection{Protocol}
\label{sec:protocol}

HBTBD results are the mean $\pm$ sample standard deviation over \textbf{ten
seeds} (five original, five fixed before any extended result was observed);
Elliptic uses the original five. Every model is re-seeded immediately before
construction, and all models within a table share an identical protocol:
features standardised on training statistics only, a strict temporal
validation split, a fixed epoch budget with best-validation checkpoint
selection, and identical loss and optimiser settings. Since every model sees
the same seeds, seeds are matched observations and we use \textbf{paired}
two-sided $t$-tests with 95\% confidence intervals. (At $n=5$ a Wilcoxon test
cannot reach $p<0.05$ at all, so rank tests are not used.) Our primary metric
is average precision (PR-AUC), the informative metric under heavy class
imbalance. We additionally report per-architecture \textbf{seed ensembles}:
the mean of the ten models' predicted probabilities, evaluated once.

This protocol supersedes the single-seed numbers in the submitted version,
and we state two corrections explicitly. First, the previously reported
69.02\% PR-AUC was an artefact of one seed in a configuration where two
models were trained in one process under one global seed; it does not
reproduce under isolated re-seeding. Second, the previous SimpleHeteroGNN
baseline was trained under a weaker protocol (random validation split,
identity metapath adjacencies), so the two were never comparable.
Table~\ref{tab:hbtbd} retrains both identically.

\subsection{Elliptic: Graph Structure and the Time Step}
\label{sec:temporal}

\begin{table}[t]
\caption{Elliptic chronological split (5 seeds). The lower block crosses
timestep position encoding with skip connections on the GraphSAGE encoder.}
\label{tab:elliptic}
\centering
\footnotesize
\setlength{\tabcolsep}{3pt}
\begin{tabular}{lcc}
\toprule
\textbf{Model} & \textbf{Test PR-AUC} & \textbf{Test ROC-AUC} \\
\midrule
Logistic Regression (no graph) & $16.21 \pm 0.00$ & $82.44 \pm 0.00$ \\
\midrule
Skip, no timestep (GraphSAGE) & $\mathbf{53.77 \pm 0.98}$ & $81.75 \pm 0.38$ \\
Skip, $+$ timestep & $52.14 \pm 1.79$ & $81.59 \pm 0.98$ \\
No skip, $+$ timestep & $43.39 \pm 4.64$ & $81.62 \pm 0.80$ \\
No skip, no timestep & $38.92 \pm 6.96$ & $79.24 \pm 1.49$ \\
\bottomrule
\end{tabular}
\end{table}

The first row of Table~\ref{tab:elliptic} isolates the value of graph
structure: moving from a feature-only classifier to GraphSAGE gains $+37.56$
PR-AUC points, 95\% CI $[+36.34, +38.78]$, $p<10^{-4}$---by far the largest
effect in the study. ROC-AUC is essentially flat across all five rows; only
average precision separates them, which is why we treat PR-AUC as primary.

The lower block crosses the timestep input with skip connections and does not
support a directional temporal claim: with skips, removing the timestep helps
by $1.63$ points ($p=0.23$); without skips, removing it \emph{hurts} by
$4.47$ points ($p=0.37$). The sign flips with the architecture and neither
contrast is significant. What \emph{is} significant is the skip connection:
$+8.75$ points with the timestep ($p=0.022$), $+14.85$ without ($p=0.007$).

\textbf{Correction and scope.} The submitted version reported that removing
temporal encoding improved PR-AUC by $6.7$ points and concluded temporal
modelling harms AML detection. That comparison rested on a single seed and on
the no-skip rows; we withdraw the claim. The variable here is only \emph{the
timestep supplied as an input feature}---we did not evaluate TGN, TGAT, or
any architecture modelling time through message passing, so nothing here
bears on temporal graph learning in general. The defensible conclusion is
methodological: a raw time index is not a reliable improvement under
chronological distribution shift, its apparent effect is easily confounded by
architectural stabilisers, and it should be ablated with seed control rather
than assumed either way.

\subsection{Heterogeneous Models on HBTBD}

\begin{table}[t]
\caption{HBTBD official split (mean $\pm$ std over 10 seeds, unified
protocol). Ens.\ rows: mean of the 10 models' probabilities, evaluated once.}
\label{tab:hbtbd}
\centering
\footnotesize
\setlength{\tabcolsep}{2.5pt}
\begin{tabular}{lccc}
\toprule
\textbf{Model} & \textbf{PR-AUC} & \textbf{ROC-AUC} & \textbf{F1} \\
\midrule
SimpleHeteroGNN     & $64.80 \pm 1.22$ & $87.62 \pm 1.20$ & $63.24 \pm 1.82$ \\
RiskMAGNN ($d$=128) & $63.50 \pm 6.78$ & $89.95 \pm 1.04$ & $62.33 \pm 3.72$ \\
RiskMAGNN ($d$=192) & $66.83 \pm 4.04$ & $90.15 \pm 1.38$ & $64.53 \pm 3.22$ \\
\midrule
SimpleHeteroGNN (ens.) & $70.18$ & $89.57$ & $67.48$ \\
RiskMAGNN $d$=192 (ens.) & $\mathbf{71.23}$ & $\mathbf{91.17}$ & $\mathbf{68.12}$ \\
\bottomrule
\end{tabular}
\end{table}

Table~\ref{tab:hbtbd} reports the central comparison. On PR-AUC no
single-model difference is significant: RiskMAGNN ($d$=192) leads
SimpleHeteroGNN by $+2.03$ points, 95\% CI $[-0.47, +4.54]$, $p=0.099$. On
ROC-AUC the advantage is significant and robust for \emph{both} sizes:
$+2.33$ points ($p=0.002$) for $d$=128 and $+2.53$ points ($p=0.002$) for
$d$=192. The honest reading is that RiskMAGNN produces a reliably better
global ranking of transactions---which is what ROC-AUC measures---but its
edge in the high-recall regime that PR-AUC rewards is within seed noise.

Two further observations matter in practice. First, \textbf{training can
collapse on individual seeds}: the $d$=128 model fell to 45.8\% PR-AUC on one
of ten seeds (hence its $\pm 6.78$), with training loss decreasing normally
throughout---only validation flags the failure. A single-seed paper hitting
such a seed (or a lucky one) would report a dramatically wrong number.
Second, the \textbf{seed ensemble} recovers what individual runs lose:
averaging the ten models' probabilities lifts every architecture by 4--5
PR-AUC points, and the RiskMAGNN $d$=192 ensemble is the best detector we
observe (71.2\% PR-AUC, 91.2\% ROC-AUC, 68.1\% F1)---the cheapest reliable
improvement available on this benchmark. The ensemble F1 narrows but does not
close the gap to the published Random Forest of Song and Gu
(73.9\%)~\cite{song2023}, whose two-stage pipeline tunes a classifier
separately on frozen features, whereas we optimise a single end-to-end
objective under focal loss, trading precision for recall on the rare class.
Selecting the F1 threshold on validation transferred \emph{worse} than the
default to test ($59.3$ vs $63.2$ for SimpleHeteroGNN)---itself evidence of
the validation--test shift documented below.

\subsection{Component Ablation}

\begin{table}[t]
\caption{Component ablation, $2\times2$ design on the RiskMAGNN $d$=128
skeleton (10 seeds). All variants share residual connections, LayerNorm,
classifier and training schedule.}
\label{tab:ablation}
\centering
\footnotesize
\setlength{\tabcolsep}{2.5pt}
\begin{tabular}{lccc}
\toprule
\textbf{Variant} & \textbf{PR-AUC} & \textbf{ROC-AUC} & \textbf{F1} \\
\midrule
Full (both)     & $\mathbf{63.50 \pm 6.78}$ & $89.95 \pm 1.04$ & $\mathbf{62.33 \pm 3.72}$ \\
$-$ TransE      & $61.54 \pm 10.18$ & $\mathbf{90.51 \pm 1.03}$ & $61.77 \pm 4.22$ \\
$-$ Risk bias   & $62.61 \pm 3.20$ & $90.46 \pm 0.47$ & $61.94 \pm 2.79$ \\
$-$ Both        & $60.77 \pm 4.04$ & $89.75 \pm 0.53$ & $60.62 \pm 3.27$ \\
\bottomrule
\end{tabular}
\end{table}

Table~\ref{tab:ablation} asks what the two proposed components contribute;
the answer is less than we initially believed, and we report the trajectory
candidly. On the original five seeds, removing both components cost $3.84$
PR-AUC points ($p=0.020$); on ten seeds the estimate is $+2.72$, 95\% CI
$[-0.10, +5.55]$, $p=0.057$---directionally consistent in nine of ten seeds
but marginal, and not claimed as confirmed. Neither component is individually
significant. What the ablation shows clearly is instability: three of four
variants---including the full model---suffer at least one collapsed seed (to
45.8\%, 34.5\% and 53.0\% PR-AUC), while uniform attention is most stable
($\pm 3.20$). Component effects of a few points cannot be reliably resolved
on a benchmark where single seeds can swing by twenty.

\subsection{The M2 Metapath Diagnostic}
\label{sec:resplit}

\begin{table}[t]
\caption{Effect of M2 availability during training (RiskMAGNN $d$=192;
official row over 10 seeds, resplit row over 5). The resplit sacrifices
temporal validity and is a diagnostic, not a performance claim.}
\label{tab:resplit}
\centering
\footnotesize
\setlength{\tabcolsep}{2.5pt}
\begin{tabular}{lccc}
\toprule
\textbf{Split} & \textbf{Val PR} & \textbf{Test PR} & \textbf{Gap (pp)} \\
\midrule
Official (M2 absent)  & $96.5 \pm 1.5$ & $66.8 \pm 4.0$ & $29.7 \pm 4.0$ \\
Resplit (M2 present)  & $97.5 \pm 1.4$ & $94.3 \pm 0.4$ & $3.2 \pm 1.4$ \\
\bottomrule
\end{tabular}
\end{table}

To isolate the effect of M2's absence we build a stratified random 80/20
resplit of the combined HBTBD data so that all three metapaths appear in
training---deliberately breaking temporal ordering, as a diagnostic only. The
validation--test gap collapses from $29.7 \pm 4.0$ to $3.2 \pm 1.4$ points
(Table~\ref{tab:resplit}), with resplit test PR-AUC remarkably stable
($\pm 0.35$): most of the gap previously attributed to model weakness is a
property of the benchmark split, which therefore understates model capacity
for layering detection and makes architecture comparisons on it comparisons
under a missing relation type.

\subsection{Explanation Findings}

Explanations for 200 test transactions show three consistent patterns:
RiskMAGNN relies mainly on neighbourhood-aggregated features (indices
41--164) rather than local attributes, validating the graph-based approach;
per-metapath edge masks assign highest importance to M1 neighbours for
illicit predictions, consistent with mixing-service behaviour; and false
positives concentrate on high-volume transactions routed through shared
output addresses---the signature of a cryptocurrency exchange. The last is
actionable: exchange-like nodes can be identified by high M1 in-degree and
given a separate threshold.

\section{Discussion}

The dominant lesson from re-running this study with seed control is about
evaluation. On HBTBD the spread across seeds ($\pm 1$ to $\pm 7$ PR-AUC
points, with individual seeds collapsing by twenty) is of the same order as
the differences between architectures, so a single-seed comparison can report
a multi-point ``improvement'' that is pure initialisation luck---as our own
submitted version did. Doubling the seeds from five to ten itself changed
which effects looked significant (the ablation effect moved from $p=0.020$ to
$p=0.057$; the ROC-AUC advantage strengthened). Multi-seed reporting with
matched-pair tests should be a baseline expectation on this benchmark.

Ordered by size, the effects that survive: graph structure ($37.6$ points),
the benchmark-split flaw ($\approx 27$ points), skip connections ($9$--$15$
points on Elliptic), seed ensembling ($4$--$5$ points for every
architecture), and RiskMAGNN's ranking advantage ($+2.3$ to $+2.5$ ROC-AUC
points, $p=0.002$). The component effect on average precision is suggestive
but unconfirmed; capacity scaling shows no significant effect. Two claims
from the submitted version did not survive and are withdrawn above. We
consider the benchmark-split flaw the most useful contribution, since it
affects anyone evaluating on HBTBD's official split.

\section{Limitations}

Ten seeds bounds seed variance but still leaves few-point effects
undetectable when single seeds can collapse, so our non-significant results
are inconclusive rather than evidence of no effect. Both benchmarks are
academic-scale (203K and 46K transaction nodes); deployment requires larger
networks with different class distributions. We do not evaluate dedicated
temporal graph architectures (TGN, TGAT), so our temporal findings are
limited to timestep position encoding. The M2 asymmetry cannot be resolved
without data from steps 1--34 in which one-way flows appear. Explanation
optimisation (200 steps per transaction) is too slow for real-time triage.
The forensic prior ordering is a design choice we initialise but do not
independently validate against labelled typologies.

\section{Conclusion}

We presented RiskMAGNN, a heterogeneous GNN for Bitcoin AML combining
translation-based metapath encoding with forensic-prior-biased attention,
evaluated over ten seeds with paired significance tests. Graph structure
dominates ($+37.6$ average-precision points over a feature-only baseline).
Our architectural contribution is real but modest: a significant ranking
advantage ($+2.5$ ROC-AUC points, $p=0.002$) whose average-precision
counterpart stays within seed noise, and a suggestive but unconfirmed joint
component effect ($+2.7$ points, $p=0.057$). The practical recommendation is
the seed ensemble, which lifts every architecture by $4$--$5$ points and
reaches 71.2\% average precision. We could not replicate the larger
single-seed gains in our submitted version and report corrected numbers here.
The most transferable result is a property of the benchmark rather than any
model: the M2 metapath is missing from HBTBD's official training split, and
that absence accounts for most of the validation--test gap. Code and all
result files are released so these numbers can be checked.

\section*{Acknowledgment}

The authors thank S P Jain School of Global Management for computational
resources and research support.

\section*{Disclosure of AI-Generated Content}

In line with IEEE policy, we disclose that an AI assistant (Anthropic's
Claude) supported this work in two ways: it helped write the experiment and
statistical-analysis scripts released with this paper, and it helped edit the
wording of the Abstract, Introduction and Discussion. All experiments were
run by the authors using the released code, and every number, claim and
conclusion was checked by the authors, who take full responsibility for the
paper's content.

\begin{thebibliography}{10}

\bibitem{song2023}
J.~Song and Y.~Gu, ``HBTBD: A heterogeneous Bitcoin transaction behavior
dataset for anti-money laundering,'' \emph{Applied Sciences}, vol.~13,
no.~15, p.~8766, Jul. 2023, doi: 10.3390/app13158766.

\bibitem{bordes2013}
A.~Bordes, N.~Usunier, A.~Garcia-Duran, J.~Weston, and O.~Yakhnenko,
``Translating embeddings for modeling multi-relational data,'' in
\emph{Advances in Neural Information Processing Systems}, vol.~26, 2013,
pp.~2787--2795.

\bibitem{fu2020}
X.~Fu, J.~Zhang, Z.~Meng, and I.~King, ``MAGNN: Metapath aggregated graph
neural network for heterogeneous graph embedding,'' in \emph{Proc. The Web
Conference 2020}, Taipei, Taiwan, Apr. 2020, pp.~2331--2341.

\bibitem{hamilton2017}
W.~L. Hamilton, R.~Ying, and J.~Leskovec, ``Inductive representation learning
on large graphs,'' in \emph{Advances in Neural Information Processing
Systems}, vol.~30, 2017.

\bibitem{weber2019}
M.~Weber \emph{et al.}, ``Anti-money laundering in Bitcoin: Experimenting with
graph convolutional networks for financial forensics,'' arXiv:1908.02591,
2019.

\bibitem{lawal2024}
I.~B. Lawal, ``An explainable graph convolutional network for anti-money
laundering detection in Bitcoin transactions,'' \emph{Int. J. Network Security
\& Its Applications}, vol.~16, no.~1, pp.~1--15, Jan. 2024.

\bibitem{ferretti2025}
F.~Ferretti, A.~Ceccarini, A.~Marchetti, and A.~Sgarro, ``A systematic
evaluation of graph neural networks for cryptocurrency anti-money
laundering,'' \emph{IEEE Access}, vol.~13, pp.~23456--23470, 2025.

\bibitem{deprez2024}
S.~Deprez, S.~Rymaszewski, and M.~Badra, ``Network analytics for anti-money
laundering: A survey of graph learning techniques,'' arXiv:2405.19383, May
2024.

\bibitem{ying2019}
R.~Ying, D.~Bourgeois, J.~You, M.~Zitnik, and J.~Leskovec, ``GNNExplainer:
Generating explanations for graph neural networks,'' in \emph{Advances in
Neural Information Processing Systems}, vol.~32, 2019, pp.~9240--9251.

\bibitem{johannessen2025}
F.~Johannessen and M.~Jullum, ``Finding money launderers using heterogeneous
graph neural networks,'' \emph{The Journal of Finance and Data Science},
vol.~11, art. no. 100175, Dec. 2025, doi: 10.1016/j.jfds.2025.100175.

\end{thebibliography}

\end{document}
